% =============================================================================
% Ethics Statement and Broader Impact
% P9 (DNB) fork of ethics_statement_zvf.tex, 2026-10-18 Deep Review fixes.
% Intended placement: after the Conclusion section, before References.
% Label: \label{sec:impact_standalone}
% Target length: >= 2 pages (NeurIPS Broader Impact guidance).
% =============================================================================

\section{Ethics, Limitations, and Broader Impact}
\label{sec:impact_standalone}

This statement is written to satisfy the NeurIPS Code of Ethics and Broader
Impact requirements. It complements the shorter Limitations and Broader
Impact subsections embedded in Section~\ref{sec:limitations} by consolidating
in one place a full dual-use analysis, itemised compute accounting, carbon
footprint estimation, data provenance disclosures, a candid acknowledgment
of the closed-source training backend on which a fraction of our headline
numbers depends, and a list of methodological limits we are aware of but
did not have resources to close within the submission window.

% ---------------------------------------------------------------------------
\subsection{Dual-Use Analysis: Misuse Risks of Reasoning RL}
\label{sec:dualuse}
% ---------------------------------------------------------------------------

\paragraph{Threat model.}
\ours{} releases (i) training scripts that apply GRPO to the GSM8K
\cite{cobbe2021training} mathematical reasoning benchmark and to the
Salesforce xLAM-Function-Calling-60k \cite{liu2024xlam} tool-use corpus,
(ii) run records (step-level reward and ZVF traces; no model weights or
LoRA adapters are released with this submission), and (iii) diagnostic code for
Zero-Variance Fraction, reward-stability, and length-bias analysis. We
analyse four concrete misuse pathways these artefacts could, in principle,
enable, and we describe the existing guardrails and the residual risk we
judge acceptable for publication.

\paragraph{M1. Weaponisation of mathematical reasoning.}
The most capability-relevant artefact we release is GRPO training code evaluated
on grade-school arithmetic and multi-step reasoning; the reviewed evidence is
inconclusive about a capability improvement. A malicious
operator could attempt to extend this pipeline to tasks of greater
dual-use concern, e.g., chemistry olympiad problems, cryptographic key
recovery, or financial fraud. We judge the \emph{marginal} uplift provided
by our release to be small for three reasons. First, the later, separate
five-seed Qwen3-8B GSM8K comparison reports $83.3\%$ versus a fixed-base
$82.0\%$; its $p{=}0.26$ ($t(4){=}1.32$) is a one-sample seed-level test against that fixed
baseline, not a paired-seed test and not reviewed-record evidence of capability
improvement. It therefore does not establish an uplift. Second, the public literature already contains GRPO implementations
\cite{shao2024deepseekmath, hu2024openrlhf, sheng2024verl, vonwerra2022trl} that are at
least as capable. Third, any reasoning uplift is bounded by the rewarder:
GSM8K rewards use exact-match against a human-written gold answer, which
does not generalise to the domains listed above without a new reward
signal, which itself requires expertise and labelled data.

\paragraph{M2. Reward hacking and the long-term alignment risk.}
GRPO, like all policy-gradient RL from a proxy reward, is vulnerable to
reward hacking: the model learns to exploit idiosyncrasies of the reward
function rather than solve the intended task \cite{skalse2022specificationgaming,
krakovna2020specification}. We observed verbosity drift in some
of our own GRPO runs, which is a concrete demonstration of this failure
mode in our own data.
Users who apply our scripts to under-specified reward functions in
production settings (e.g., a custom ``helpfulness'' classifier) may observe
models that appear to improve on held-out metrics while silently
optimising against unintended proxies. We include the Zero-Variance
Fraction and reward-stability diagnostics precisely so that downstream
users can inspect reward telemetry; detection of policy drift or reward hacking
requires separate validation. ZVF measures within-group reward contrast.

\paragraph{M3. Circumventing safety training via distillation or tool-use.}
The tool-use scripts train LoRA adapters intended to teach a base model to emit
schema-valid JSON function calls \cite{liu2024xlam}; the two tool-use case
studies in the manifest reached zero reward and no tool-use adapter is released. A motivated adversary
could in principle use the same pipeline to teach a base model to emit
jailbreak prompts as ``tool calls'' or to invoke an adversarial external
tool. We mitigate this by (a) releasing no trained weights or adapters,
and (b) not releasing any function-calling harness that would actually
execute emitted calls.

\paragraph{M4. Compute-efficiency externalities.}
Our case studies run critic-free RL on modest budgets. Improved training efficiency is
itself dual-use: it lowers the cost of running adversarial fine-tuning at
scale. We judge this externality to be modest relative to the already-low
cost of fine-tuning a sub-10B model with commodity LoRA recipes
\cite{hu2022lora}, but we flag it explicitly.

\paragraph{Residual risk.}
After weighing M1--M4 we judge that the reproducibility, calibration, and
pedagogical benefits of releasing \ours{} outweigh the residual misuse
risk. We release no trained weights, and we publish reward diagnostics without
claiming validated reward-hacking detection.

% ---------------------------------------------------------------------------
\subsection{Compute Cost Accounting}
\label{sec:compute_accounting}
% ---------------------------------------------------------------------------

The project did not retain a per-run billing ledger or measured
GPU-hour log (no such file exists in the released repository), so earlier
drafts' itemised spend and GPU-hour tables were estimates without a source
artifact and have been removed. We report compute only qualitatively:
Tinker runs used a serverless commercial API; open-stack runs used Modal
H100 and L4 workers; failed runs consumed compute too and are listed in
Table~\ref{tab:run-manifest}.

\paragraph{Hidden costs.} Beyond direct platform spend, the project
consumed human labour (approximately six student-FTE-months, unpaid) and
infrastructure subsidised by an academic research group with
shared access to one A100 node used for TRL baselines and
pre-submission sanity checks. No author received commercial compensation
from Thinking Machines, Modal, or any other vendor mentioned.

% ---------------------------------------------------------------------------
\subsection{Carbon Footprint Estimation}
\label{sec:carbon}
% ---------------------------------------------------------------------------

Because GPU-hours were not logged per run (Section~\ref{sec:compute_accounting})
and Tinker does not disclose hardware or energy telemetry, we do not report a
numerical energy or CO$_2$-equivalent estimate; any figure would rest on
assumed rather than measured GPU-hours.

\paragraph{Offset and mitigation.}
We did not purchase carbon offsets. The run manifest
(Table~\ref{tab:run-manifest}) records which cells failed or were partial,
so that replicators can skip them, and marks the cells for which step-level
logs and checkpoints are retained. We regard reproducibility-by-artefact as a
more durable mitigation than offsets.

% ---------------------------------------------------------------------------
\subsection{Data Provenance}
\label{sec:provenance}
% ---------------------------------------------------------------------------

\ours{} uses only publicly released research datasets. No private data,
personally identifiable information, or licensed proprietary content is
used. No human annotators were employed during this work. We document
each dataset below.

\paragraph{GSM8K \cite{cobbe2021training}.}
8{,}792 grade-school math word problems (7{,}473 train / 1{,}319 test)
authored by human writers contracted by OpenAI. Released under the
\textbf{MIT License} via \url{https://github.com/openai/grade-school-math}.
We draw training prompts from the standard train split and held-out
evaluations from the test split without modification (individual held-out
evaluations use test-split subsets whose sizes are stated where they are
reported; none of them is the full 1{,}319-problem test set unless stated). The canonical
citation is \citet{cobbe2021training}. Known limitations of GSM8K include
gender-neutral but culturally US-centric word problems (names, currencies,
sports) and a moderate rate of human-labelling errors estimated at
\textasciitilde 2\% by \cite{zhang2024gsm1k}. Our held-out evaluation
respects the test/train split.

\paragraph{Salesforce xLAM-Function-Calling-60k \cite{liu2024xlam}.}
60{,}000 function-calling examples generated by Salesforce Research as
training data for the xLAM agentic model family. Released under the
\textbf{CC BY 4.0} license via \url{https://huggingface.co/datasets/Salesforce/xlam-function-calling-60k}.
The training scripts accept the public release as-is; no xLAM-trained run
is reported in this manuscript (the two tool-use case studies in
Table~\ref{tab:run-manifest} use the authored synthetic corpus below).
Known limitations: xLAM schemas are synthetic and skew toward cleanly-typed
arguments; the dataset under-represents ambiguous, error-handling, and
multi-turn tool-calling. The dataset is by Salesforce Research \cite{liu2024xlam}, credited here as
the CC BY 4.0 license requires.

\paragraph{HumanEval \cite{chen2021codex}.}
164 Python programming problems with unit tests, released by OpenAI under
the \textbf{MIT License}. Its single cell in the manifest failed before
training (Tier C), so no HumanEval result is reported. Known
limitations: small scale, English-language docstrings, single-language
coverage; documented contamination concerns against frontier pretraining
corpora \cite{riddell2024contamination}.

\paragraph{NuminaMath \cite{li2024numinamath}.}
Named in an external collaborator's planned multi-stage
GSM8K+NuminaMath pipeline (Apache 2.0). Those collaborator cells never
started (Tier C in Table~\ref{tab:run-manifest}); no derivative weights are
released.

\paragraph{Open-Platypus \cite{lee2023platypus}.}
Named in an external collaborator's planned Qwen3-8B code-generation
warm-up (curated subset, \textbf{CC BY-NC 4.0}). That cell never started
(Tier C); nothing derived from it is released.

\paragraph{Synthetic tool-use corpus (authored).}
The two tool-use case studies use ten hand-written prompts over five
stub tools (calculator, weather, search, reminder, translate), defined
inline in \texttt{experiments/tinker-runs/scripts/tinker\_parallel\_runner.py}
(\texttt{get\_tool\_use}). The shaped reward (0.3 valid JSON, $+0.4$ correct
tool, $+0.3$ required arguments) parses the \emph{entire} response as JSON,
and the prompt uses the Qwen chat template for every model, so the zero
reward of both runs may reflect output formatting rather than tool-use
ability. The prompts are authored by the paper's authors and contain no
third-party content, user data, or proprietary API traces.

\paragraph{No web scraping, no human subjects.}
We did not scrape any website. We did not run any study with human
participants. No IRB review was therefore required. No data that could
reasonably be considered to contain PII, copyrighted content, or sensitive
personal attributes was used at any stage of training, evaluation, or
reward modelling.

% ---------------------------------------------------------------------------
\subsection{Closed-Source Tinker Acknowledgment}
\label{sec:tinker_closed}
% ---------------------------------------------------------------------------

A central tension in \ours{} is that a substantial fraction of our
headline numbers come from a closed-source commercial platform while our
paper simultaneously advocates for reproducibility and platform
independence. We do not resolve this tension by hiding it; we describe
it precisely.

\paragraph{What Tinker is and is not.}
Tinker \cite{thinkingmachines2024tinker} is a managed LLM fine-tuning and
inference service provided by Thinking Machines, Inc. It exposes a Python
SDK (\texttt{tinker==0.16.1}) that accepts custom loss functions
(\texttt{forward\_backward\_custom}), a limited set of optimisers, and
standard LoRA hyperparameters. It does not expose (i) the exact server-side
GRPO loss implementation, (ii) reward normalisation or baseline subtraction
scheme, (iii) minibatch construction or gradient-accumulation strategy,
(iv) hardware configuration (GPU type, inter-node bandwidth), or (v)
system-level telemetry (energy, throughput, queueing).

\paragraph{What this means for our claims.}
Tinker results in this paper measure the \emph{platform's} implementation
of GRPO, not an abstract specification of the algorithm. Cross-stack reward
gaps between Tinker GRPO and an open-source GRPO implementation cannot be
fully attributed from our data: we are able to rule out a handful of
candidate explanations (seed variance, model-size confound) but not the
implementation itself. We therefore treat only the open-stack Modal cells (TRL and the
classic-RL libraries; no veRL or OpenRLHF run is in the manifest) as
hyperparameter-auditable, and use Tinker results as
\emph{descriptive case studies} of what a carefully engineered production
stack can achieve for critic-free RL at our scales.

\paragraph{Reproducibility commitments we can make.}
\begin{enumerate}
  \item All Tinker experiment scripts, configuration files, JSON step
        logs, and per-run W\&B projects are archived in the repository.
        Researchers with Tinker access can attempt replication with our
        exact configurations.
  \item Every table row and summary statistic derived from Tinker data is
        labelled as such by its stack column (Table~\ref{tab:artifact-card}),
        backend column (Table~\ref{tab:run-manifest}), or block header
        (Table~\ref{tab:main_results}); independent replication requires
        Tinker API access.
  \item The paper makes no empirical inferential claim; the centered
        reward-contrast identity (Appendix~\ref{sec:appendix:zvf-formalization})
        is a mathematical statement under the stated runner assumptions.
        The Tinker-derived stratified ZVF
        correlations (Section~\ref{sec:zvf-analysis}) are reported as
        artifact analyses with explicit closed-backend caveats and
        single-seed-per-cell limitations, not as platform-independent
        statistical conclusions.
  \item We commit to re-running any Tinker-only experiment on an open
        backend if an equivalent open-source service becomes available,
        and to issuing a revised version of this paper if any Tinker
        result is ever shown to be due to an undisclosed
        implementation choice rather than the algorithm.
\end{enumerate}

\paragraph{Key rotation and credential hygiene.}
Our Tinker API key (prefix \texttt{tml-...}) is held in a repository
\texttt{.env.example} placeholder only. The real key is stored in the
authors' password manager and rotated whenever a failure mode suggests
possible exfiltration. No Tinker key is committed to git history in the
anonymous source repository
(\url{https://anonymous.4open.science/r/agentic-grpo-bench-anon-2C6B}).

% ---------------------------------------------------------------------------
\subsection{Known Methodological Limits}
\label{sec:methodlimits}
% ---------------------------------------------------------------------------

In addition to the infrastructure failures and platform-specific caveats
enumerated in Section~\ref{sec:limitations}, we flag the following
methodological limits.

\paragraph{Short training horizons.}
All Tinker GRPO runs in the manifest were capped at 10--30 gradient steps
(Table~\ref{tab:run-manifest}). This is a deliberate cost-control choice and
means our Tinker numbers are \emph{early-training snapshots}. We cannot rule
out that longer training would change the qualitative picture.

\paragraph{Single-seed Tinker experiments.}
Each Tinker configuration was run once, not 3--5 times as best practice
prescribes \cite{henderson2018deep}. Tinker results therefore lack
variance estimates and do not support significance testing; we report
them descriptively.

\paragraph{Train-set reward as primary Tinker metric.}
Tinker runs primarily report reward on training prompts. Held-out
evaluation exists only for the cells marked in the Held-out column of
Table~\ref{tab:run-manifest}; for all other cells we cannot distinguish
memorisation from generalisation.

\paragraph{LoRA only; no full fine-tuning.}
All language-model experiments use LoRA \cite{hu2022lora}. We have
\emph{not} tested whether full fine-tuning would re-order the
library/algorithm comparisons. The LoRA constraint is practical (cost) but
it means our claims about PPO vs.\ GRPO and TRL vs.\ Tinker are
LoRA-conditional.

\paragraph{Benchmark coverage is narrow.}
Completed cells cover synthetic arithmetic, GSM8K, and a synthetic
tool-use corpus; the single HumanEval cell failed (Tier C). We do not evaluate on MT-Bench,
ArenaHard, safety benchmarks (HarmBench, ToxicChat), or truthfulness
benchmarks (TruthfulQA). Any claim about ``reasoning'' in the paper is
strictly a claim about the covered benchmarks.

\paragraph{No human preference data.}
We use verifiable rewards only (exact-match for math, unit-test for code,
schema-validity for tool-use). We do not study RLHF with human preference
data; findings should not be extrapolated to reward-model--based alignment
without further work.

\paragraph{Closed-source implementation opacity (reiterated).}
Comparisons between Tinker GRPO and TRL GRPO are confounded by the
closed-source nature of Tinker's implementation. See
Section~\ref{sec:tinker_closed} for detail.

\paragraph{Cross-platform hardware confounding.}
Tinker, Modal H100, and Modal L4 use different accelerators,
different memory hierarchies, and different inference/training stacks
(Tinker proprietary; TRL and the classic-RL libraries on Modal).
Observed differences in reward dynamics are not purely algorithmic.

\paragraph{Exploratory, not confirmatory.}
This is an exploratory case study, not a pre-registered confirmatory
experiment. We did not pre-register primary hypotheses. All $p$-values are
un-corrected for the number of comparisons actually made across the
paper, and all of our analyses are descriptive.

\paragraph{Geographic and demographic limits.}
The specific city and country of the authors' institution are withheld
here for anonymity but disclosed in the non-anonymous version of the
paper. Our choice of benchmarks, prompt phrasings, and evaluation design
reflects that single-region context. Independent replication by groups
in other regions, on non-English tasks, and with non-Qwen / non-Llama
families is an open question.

% =============================================================================
% End of ethics_statement.tex
